% ============================================================
% 朗读者·名言控（杂志页码 42–45）
% ============================================================
\colchapter{朗读者·名言控}

\section*{李玉刚：“以传统文化向世界传递大爱”}

\auth{《开讲啦》演讲摘录}

大家好！我是李玉刚，一名歌者，同时是一名舞台剧演员。

我首先想给大家讲一个故事。一个小男孩出生在吉林省公主岭市，他在家里排行老三，乳名叫作“二弟”，他就这样不起眼地生长在中国东北的农村。小时候，他跟所有农村孩子一样，特别淘气，但这个小男孩还有一个显著的特点——会唱歌，那时候，他的歌声总是飘荡在东北的黑土地上。那时，物质生活非常匮乏，在他十岁左右，邻居家买了村里的第一台电视机，他有幸看到电视，而第一次看到的电视节目就是春节联欢晚会。没想到，若干年之后，这个不起眼的小孩也登上了春晚，演唱了一首《新贵妃醉酒》。讲到这里，大家可能猜到了，这个小男孩就是我。

高中毕业后，我考入吉林艺术学院，但我还是选择了放弃，于是，18岁的我从公主岭走向了社会。我从一个餐厅的服务员做起，餐厅的中心有一个舞台，客人可以边吃饭边听演员唱歌，后来，我自己也站上了那个舞台。再后来，我离开餐厅也离开东北，去了西安，还有幸遇到了我的京剧启蒙老师，马洪才老师。

我决定学习京剧后，为了在舞台上有更好的综合表演，又开始学习形体、舞蹈以及京剧里边的手眼身法，而那需要压腿。现在回想起20岁时压腿的情境，仍然不寒而栗，感觉就像上刑一样，但还好我咬牙坚持了下来。

到现在，我都记得老师的教诲：不要当明星，要在被别人熟知的基础上，传达一种价值观、一种艺术理念。我问他价值观、艺术理念从哪里来，他只简单地说了一句：“一切都从中国5000年悠久的文化历史当中来。”当时我也似懂非懂，这简单的一句话，我可以用一生去解读它。登上春晚之后，我一直在想自己究竟该向哪方面发展，就做一名所谓的歌手，当然也不错，但我对自己还是有了更高的要求——我希望自己是一名戏剧演员，也希望自己是一名戏剧导演。

之后，我导演并主演了自己人生当中非常重要的一部舞台作品——歌舞剧《昭君出塞》。那是2013年，为了找到创作灵感，我和团队重走了王昭君的路，从她的家乡一路向北。快到呼和浩特的一个黄昏，残阳如血，孤独感一下涌上了心头，可心中的信念却不能放下。那一刻，我忽然体会到了王昭君内心的复杂与坚定：她离开故土，走向大漠，将中原的文明带向远方，用一生换来两个民族的和平。以传统文化向世界传递大爱，这也是我艺术生涯中一直坚守的方向。

在我的歌词中，有取之不尽的中国美学意象；在我的舞台上，有穿越千年的东方气韵。中华文化是人类文明的瑰宝，其中有取之不尽的精神财富。作为艺术工作者，我们应该让世界看到中国之美，让更多人因为我们的作品而爱上这片土地。谢谢大家！

\begin{quoteBox}{赏句录}
坚持的能量也许并不会瞬间惊天动地，但会让这个世界悄然发生变化。正如一滴水的力量虽然弱小，但经年累月的坚持，也能让岩石败下阵来。滴水穿石，依靠的不是强大的力量，而是“纵使疾风起，人生不言弃”的处世之态。
\end{quoteBox}

\begin{quoteBox}{赏句录}
艺术的生命力在于扎根传统的土壤，又在时代的雨露中长出新的枝叶。唯有心怀敬畏、笃行不怠，方能让古老的文明在当代舞台上焕发出夺目光彩，也唯有如此，才能让文化自信成为我们内心最深沉、最持久的力量。
\end{quoteBox}

\clearpage
\section*{徐默凡：“锦绣文章也永远事关人间烟火”}

\auth{华东师范大学中文系开学典礼致辞摘录}

各位同学，大家好。今天，我想和大家谈一谈关于写作和语言的三件事，或者说是三个“陷阱”。

第一个陷阱是“强行分类”。进入大学之后，我们接触的专业知识越来越系统，分类的意识也越来越强。但语言和文学的世界里，很多美好的东西恰恰生长在分类的缝隙之中。如果遇到任何作品都先想“这是现实主义还是浪漫主义”“这是抒情散文还是议论文”，反而会遮蔽我们最直接的感受。分类是理解的工具，而不是感受的牢笼。

第二个陷阱是“隔膜体验”。在互联网时代，我们习惯于通过屏幕去认识世界，很多体验是二手甚至三手的。文字是符号的符号，如果我们的生活体验越来越稀薄，那么写出来的锦绣文章，也终将失去根基。卡西尔说“人是符号的动物”，但符号的源头永远是鲜活的、有温度的生活。

第三个陷阱是“以形害意”。我们追求辞藻的华美，却常常忘了语言的根本目的是传情达意。有人问，什么样的文章是好文章？我以为，好的文章不是堆砌了多少漂亮话，而是有没有说清楚一件事、有没有打动人。锦绣文章也永远事关人间烟火——真正动人的文字，永远来自对生活的热爱与真诚的观察。

写作从来不是技巧的堆砌，而是一种与世界对话的方式。希望你们在未来的日子里，多读、多想、多写，让文字成为你们认识世界、表达自我的桥梁。

\begin{quoteBox}{赏句录}
分类是理解的工具，而不是感受的牢笼。当我们习惯用一个个现成的标签去框定作品与人生时，恰恰会错过那些正在缝隙里蓬勃生长的惊喜。
\end{quoteBox}

\begin{quoteBox}{赏句录}
人是符号的动物，但符号的源头永远是鲜活的、有温度的生活。锦绣文章也永远事关人间烟火——真正动人的文字，永远来自对生活的热爱与真诚的观察。
\end{quoteBox}

\begin{quoteBox}{写作点拨}
“三个陷阱”的演讲结构值得借鉴：以分类的方式将复杂话题拆解，先亮观点，再举例分析，最后总结升华。这种“总分总”的层次让说理清晰而有条理，适合迁移到议论文的结构搭建中。
\end{quoteBox}
